% Intended LaTeX compiler: xelatex
% Layout modeled on a two-column academic CV: section labels in a left
% column (light sans small caps), content on the right, Garamond body text.
\documentclass[11pt]{article}
\usepackage{fontspec}
\usepackage[left=1.9in,right=0.7in,top=0.6in,bottom=0.8in]{geometry}
\usepackage[explicit]{titlesec}
\usepackage{changepage}
\usepackage{enumitem}
\usepackage[hidelinks]{hyperref}
\urlstyle{same}

% Fonts: EB Garamond (body, stands in for Sabon) and Alegreya Sans SC Light
% (section labels, stands in for Cronos Pro Light). Both ship with TeX Live.
\setmainfont{EBGaramond}[
  Extension      = .otf,
  UprightFont    = *-Regular,
  ItalicFont     = *-Italic,
  BoldFont       = *-Bold,
  BoldItalicFont = *-BoldItalic,
  Numbers        = Lining,
  Scale          = 1.08]
\newfontfamily\labelfont{AlegreyaSansSC-Light.otf}

% Width of the left label column (it hangs into the left margin)
\newlength{\cvlabelwidth}
\setlength{\cvlabelwidth}{1.35in}

% Section labels sit in the left margin, top-aligned with the first line
\titleformat{\section}[leftmargin]
  {\labelfont\large\raggedright}{}{0pt}{#1}
\titlespacing*{\section}{\cvlabelwidth}{10pt}{0pt}

\setlength{\parindent}{0pt}
\setlength{\parskip}{0pt}
\setlength{\emergencystretch}{2em}
\pagestyle{empty}

% Dated entry: bold first line with the date flush right, remaining lines indented
%   \entry{date}{first line}{other lines, separated by \\}
\newcommand{\entry}[3]{%
  \textbf{#2}\hfill #1\par\nopagebreak
  \begin{adjustwidth}{1.5em}{0pt}\raggedright #3\end{adjustwidth}\vspace{7pt}}

% Numbered paper, counted continuously across sections:
%   \paper{bold title}{coauthors (italic)}{status line, may be empty}
\newcounter{paper}
\newcommand{\paper}[3]{%
  \stepcounter{paper}%
  {\raggedright\textbf{\thepaper. #1}\par}\nopagebreak
  \emph{#2}\par
  \if\relax\detokenize{#3}\relax\else #3\par\fi
  \vspace{7pt}}

\hypersetup{
 pdfauthor={Rafael Campello de Alcantara},
 pdftitle={Curriculum Vitae},
 pdflang={English}}

\begin{document}

% Header: name and title on the left, contact on the right, full-width rule
\begin{adjustwidth}{-\cvlabelwidth}{0pt}
\begin{tabular}[b]{@{}l@{}}
  {\fontsize{18}{22}\selectfont\scshape Rafael Campello de Alcantara}\\[2pt]
  {\itshape Curriculum Vitae}
\end{tabular}%
\hfill
{\addfontfeatures{Numbers=OldStyle}%
\begin{tabular}[b]{@{}r@{}}
  The University of Texas at Austin\\
  \href{mailto:rafael.campellodealcantara@mccombs.utexas.edu}{rafael.campellodealcantara@mccombs.utexas.edu}\\
  \url{https://rafaelcalcantara.github.io/}
\end{tabular}}

\vspace{6pt}
\noindent\rule{\linewidth}{0.5pt}
\end{adjustwidth}
\vspace{4pt}

\section{Research experience}
\label{sec:research.experience}
\entry{2024-present}{Postdoctoral fellow}{%
  The University of Texas at Austin\\
  McCombs School of Business}
\entry{2023}{Visiting scholar}{%
  School of Mathematical and Statistical Sciences, Arizona State University\\
  Sponsor: Paul Richard Hahn}
\entry{2022}{Insper, São Paulo, Brazil}{%
  Research Assistant\\
  Supervisors: Naércio Aquino Menezes Filho, Camila de Freitas Souza Campos}

\section{Education}
\label{sec:org7da8f15}
\entry{2020-2024}{PhD in Business Economics}{%
  Insper - Institute of Education and Research, São Paulo, Brazil\\
  Dissertation: A Constrained BART Model For Identifying Heterogeneous Treatment Effects In Regression Discontinuity Designs\\
  Advisor: Hedibert Freitas Lopes}
\entry{2018-2020}{MSc in Applied Economics}{%
  Universidade de São Paulo, USP, São Paulo, Brazil\\
  Thesis: Estimation of the R\&D Elasticity in Brazil: An Unobserved Common Factor Approach\\
  Advisor: Sérgio Kannebley Júnior}
\entry{2013-2017}{BA in Economics}{%
  Universidade de Brasília, UnB, Brasília, Brazil\\
  Advisor: Andrea Felippe Cabello}

\section{Publications}
\label{sec:publications}
\paper{\href{https://arxiv.org/abs/2503.00326}{A Bayesian Additive Regression Tree Model for Learning Conditional Average Treatment Effects in Regression Discontinuity Designs}}
  {with P. Richard Hahn and Hedibert F. Lopes}
  {Accepted for publication, Observational Studies}
\paper{\href{https://revistas.usp.br/ecoa/pt_BR/article/view/190655/214830}{R\&D Elasticity in Brazilian Manufacturing and Spillovers - An Unobserved Common Factors Approach}}
  {with Sergio Kannebley Jr.}
  {Brazilian Journal of Applied Economics, 2024}

\section{Preprints\slash Submitted}
\label{sec:preprints}
\paper{\href{https://arxiv.org/abs/2407.14365}{Modified BART for Learning Heterogeneous Effects in Regression Discontinuity Designs}}
  {with Meijia Wang, P. Richard Hahn and Hedibert F. Lopes}
  {}

\section{Work in progress}
\label{sec:working.papers}
\paper{Searching for Parallel Trends: Discovering Diff-in-Diff Estimators,}
  {with Andrew Herren, Palak Jain, P. Richard Hahn, Jared Murray}
  {}
\paper{\href{https://arxiv.org/abs/2208.09970}{Statistical Aspects of SHAP: Functional ANOVA for Model Interpretation}}
  {with Andrew Herren, P. Richard Hahn}
  {}

\section{Seminar and conference presentations}
\label{sec:presentations}
\begin{itemize}[label={}, leftmargin=0pt, itemsep=7pt, topsep=0pt]
  \item \emph{Searching for parallel trends: Discovering diff-in-diff estimators} (with Andrew Herren, Palak Jain, P. Richard Hahn, Jared Murray), European Seminar on Bayesian Econometrics (ESOBE), August 2026
  \item \emph{Learning Conditional Average Treatment Effects in Regression Discontinuity Designs using Bayesian Additive Regression Trees} (with P. Richard Hahn, Hedibert F. Lopes), International Society for Bayesian Analysis (ISBA) World Meeting, July 2026
  \item \emph{A Constrained BART Model for Identifying Heterogeneous Treatment Effects in regression discontinuity design} (with Meijia Wang, P. Richard Hahn, Hedibert F. Lopes), LACEA-LAMES Annual Meeting, November 2024
  \item \emph{A Modified BART Prior for Regression Discontinuity Designs} (with Meijia Wang, P. Richard Hahn, Hedibert F. Lopes), NBER-NSF Seminar on Bayesian Inference in Econometrics and Statistics, August 2023
  \item \emph{Bayesian Additive Regression Trees For Regression Discontinuity Designs} (with Hedibert F. Lopes), 44th Meeting of the Brazilian Econometric Society, December 2022
\end{itemize}

\section{Teaching}
\label{sec:org3c7d8d1}
\entry{2023}{Insper, São Paulo, Brazil}{%
  Teaching assistant of Bayesian Econometrics\\
  Professor: Hedibert Freitas Lopes}
\entry{2022}{Insper, São Paulo, Brazil}{%
  Teaching assistant of Bayesian Learning\\
  Professor: Hedibert Freitas Lopes}
\entry{2019}{Universidade de São Paulo, USP, São Paulo, Brazil}{%
  Teaching assistant of Macroeconomics 1\\
  Professor: Rudinei Toneto Júnior}
\entry{2019}{Universidade de São Paulo, USP, São Paulo, Brazil}{%
  Teaching assistant of Introduction to Macroeconomics\\
  Professor: Rudinei Toneto Júnior}
\entry{2015}{Universidade de Brasília, UnB, Brasília, Brasil}{%
  Teaching assistant of Economic Formation of Brazil\\
  Professor: Flávio Rabelo Versiani}

\section{Technical skills}
\label{sec:orgc51371c}
Advanced R, intermediate Stata, basic SQL, basic Python

\section{Language}
\label{sec:orgf2accdc}
Native Portuguese, fluent English
\end{document}
